\section{Printer Lights and the Color Science Engine}
\label{sec:printerlights}

\subsection{The Printer Light Unit}

In an optical printer, each of three channels of the printing beam is attenuated
by a calibrated filter wheel. The industry-standard increment --- the
\emph{printer point} --- is

\begin{equation}
1 \text{ point} \;=\; 0.025 \text{ in } \log_{10} \text{ exposure},
\end{equation}

so that exactly twelve points make one stop, since $12 \times 0.025 = 0.30
\approx \log_{10} 2$. Printer lights are specified as an integer triple, and a
grade is communicated between laboratory and client as three numbers such as
$(28, 25, 26)$.

\EM{} adopts this unit exactly, as an offset from the computed aim balance of
\eqref{eq:aimbalance}:

\begin{equation}
\boxed{\;\log_{10} L_c \;=\; \log_{10} L_c^{\mathrm{aim}} \;+\; 0.025\, p_c\;}
\label{eq:printerlights}
\end{equation}

with $p_c \in [-12, +12]$ integer, giving a $\pm 1$ stop range per channel. The
control is deliberately \emph{integer-valued}: printer points are integers in
practice, the quantization is finer than the visual threshold, and integer
values are what makes a grade communicable. A continuous slider here would be a
small betrayal of parameter honesty for no benefit.

\subsection{Sign Convention}

Increasing a printer light increases the exposure of the print stock in that
record, which develops more of the complementary dye, which \emph{subtracts}
that color from the print. Raising the red printer light makes the print
\emph{more cyan}, and raising all three makes the print darker.

This is counterintuitive to anyone who has used a digital color balance control,
and it is correct. \EM{} defaults to the laboratory convention and labels it
unambiguously, with the direction indicated on the control itself
(Section~\ref{sec:ui}). A preference exists to invert the sign for users who
prefer the digital convention; the underlying parameter is unchanged and
recipes remain interchangeable.

\subsection{Why Printer Lights Are Not RGB Gain}

This is the central claim of the section and the justification for exposing the
control at all.

An RGB gain applied to a display-referred image multiplies the output. A printer
light offsets the print stock's \emph{input} in log exposure, before the print
characteristic curve. Composing \eqref{eq:printerlights} with
\eqref{eq:printcurve}, the change in print density resulting from a small
printer light change is

\begin{equation}
\frac{\partial D'_c}{\partial p_c} \;=\; 0.025\;\gamma'_c(\logE'_c),
\label{eq:plsensitivity}
\end{equation}

where $\gamma'_c(\cdot)$ is the print stock's \emph{point} gamma from
\eqref{eq:pointgamma}. The sensitivity is therefore proportional to the local
slope of the print curve: large in the mid-scale, and vanishing in both the toe
and the shoulder.

The consequences are exactly the behaviors that photographers describe as
"filmic":

\begin{itemize}
\item \textbf{Color balance changes do not tint the highlights.} Near the
print's toe (scene highlights), $\gamma' \to 0$, so
$\partial D'/\partial p \to 0$. Whites stay white no matter how far the balance
is pushed. An RGB gain tints them immediately.
\item \textbf{Color balance changes do not tint the deep blacks.} The same
argument at the shoulder.
\item \textbf{The effect is strongest exactly where the eye judges color} ---
the mid-scale, skin tones, foliage.
\end{itemize}

Figure~\ref{fig:plresponse} shows the sensitivity curve. Reproducing this shape
is impossible with any pointwise gain applied after tone mapping, which is the
practical reason a print-domain model earns its complexity.

\begin{figure}[htbp]
\centering
\begin{tikzpicture}
\begin{axis}[
  width=\columnwidth, height=4.2cm,
  xlabel={scene $\logE$ (relative)},
  ylabel={$\partial D' / \partial p$},
  xmin=-2.6, xmax=1.4, ymin=0, ymax=0.09,
  grid=both, grid style={draw=emLine!20},
  label style={font=\scriptsize}, tick label style={font=\scriptsize},
  samples=250,
]
\addplot[emKey, very thick, domain=-2.6:1.4]
  {0.025*2.80*( 1/(1+exp(-(2.80*((1.05 - 0.61*(x+2.10)) + 0.55))/0.07))
              - 1/(1+exp(-(2.80*((1.05 - 0.61*(x+2.10)) + 0.55) - 2.45)/0.06)) )};
\end{axis}
\end{tikzpicture}
\caption{Printer light sensitivity \eqref{eq:plsensitivity} as a function of scene exposure. Color balance authority is concentrated in the mid-scale and vanishes at both extremes. This self-protecting behavior is a property of the model, not a safeguard added on top of it.}
\label{fig:plresponse}
\end{figure}

\subsection{Density (Overall Print Exposure)}

The print exposure time $t_{\mathrm{print}}$ of \eqref{eq:printexposure} is the
lightness control, and is exposed as \emph{print density} in the same printer
point unit applied to all three channels simultaneously:

\begin{equation}
\log_{10} t_{\mathrm{print}} = 0.025\, p_{\mathrm{master}}, \quad p_{\mathrm{master}} \in [-24, +24].
\end{equation}

Increasing density darkens the print. This control replaces "brightness" and
"exposure" in the develop workspace; the capture-time exposure compensation of
\eqref{eq:anchor} remains separate and available, and the distinction between
them is real and worth teaching: exposure compensation changes where the image
sits on the \emph{negative}'s curve (affecting shadow detail and grain), while
print density changes where it sits on the \emph{print}'s curve (affecting only
the print rendering). Users discover this because the two controls genuinely
behave differently.

\subsection{Additional Color Science Controls}

Beyond printer lights, four controls operate in the print domain. Each is
defined mathematically here and mapped to its UI representation in
Table~\ref{tab:controlmap}.

\paragraph{Saturation density} Rather than a chroma multiplier, saturation is
controlled by scaling the off-diagonal terms of the printing density matrix
$\mathbf{C}$ of \eqref{eq:cmatrix}:

\begin{equation}
\mathbf{C}(\varsigma) = \mathbf{I} + \varsigma\,(\mathbf{C} - \mathbf{I}) + (1-\varsigma)\,\varpi\,\mathbf{J},
\label{eq:satdensity}
\end{equation}

where $\mathbf{J}$ is the matrix of ones scaled to preserve neutrals and
$\varsigma$ is the control. Reducing $\varsigma$ reduces the unwanted-absorption
crosstalk, \emph{increasing} apparent saturation; raising it above unity adds
crosstalk and desaturates. Because the operator lives before the print curve, its
effect is again density-dependent: saturated shadows and saturated highlights
behave differently, which a chroma multiplier cannot express.

\paragraph{Highlight roll-off} Directly modifies the print stock's toe softness
$\kappa'_t$ over $[0.5, 2.0]\times$ nominal. This is the honest control for what
other applications call "highlights."

\paragraph{Shadow lift} Modifies $\Dmax'$, which is the print's maximum density
and therefore its black level. Reducing $\Dmax'$ produces the lifted,
low-contrast black of a print made on aged paper or a flashed negative.

\paragraph{Neutral axis} A pair of parameters $(\delta_{RG}, \delta_{BG})$
applying a density-dependent tilt to the neutral axis:

\begin{equation}
D'_R \mathrel{+}= \delta_{RG}\, \psi(D'), \qquad
D'_B \mathrel{+}= \delta_{BG}\, \psi(D'),
\label{eq:neutralaxis}
\end{equation}

with $\psi(D') = (D' - \Dmin')/\Delta D' - \tfrac{1}{2}$ an antisymmetric
weighting. This is split toning expressed as an axis tilt: positive
$\delta_{RG}$ warms shadows and cools highlights simultaneously, exactly as a
tilted neutral axis in a real print does, and unlike independent
shadow/highlight tint controls it cannot produce a physically impossible
non-monotone neutral.

\subsection{Control Mapping Table}

Table~\ref{tab:controlmap} is normative. It is the authoritative translation
between the user-visible vocabulary, the internal parameter, and the equation
that consumes it. Any control added to the product must be added to this table
or it does not ship.

\begin{table}[htbp]
\caption{Control Mapping: User Term to Model Parameter}
\label{tab:controlmap}
\begin{center}
\renewcommand{\arraystretch}{1.12}
\begin{tabularx}{\columnwidth}{@{}>{\raggedright\arraybackslash}X l l@{}}
\toprule
\textbf{User control} & \textbf{Parameter} & \textbf{Eq.} \\
\midrule
Exposure compensation & $\varepsilon$ & \eqref{eq:anchor} \\
Film speed / rating & $S$ & \eqref{eq:isoshift} \\
Push / pull & $n \to A$ & \eqref{eq:activity} \\
Development time & $t/t_0$ & \eqref{eq:activity} \\
Development temperature & $T$ & \eqref{eq:activity} \\
Agitation & $a_g$ & \eqref{eq:activity} \\
Contrast index (readout) & CI & \eqref{eq:ci} \\
Coupler activity & $\|\mathbf{A}\|$ & \eqref{eq:interlayer} \\
Printer lights R/G/B & $p_R, p_G, p_B$ & \eqref{eq:printerlights} \\
Print density & $p_{\mathrm{master}}$ & \eqref{eq:printexposure} \\
Saturation density & $\varsigma$ & \eqref{eq:satdensity} \\
Highlight roll-off & $\kappa'_t$ & \eqref{eq:printcurve} \\
Shadow lift & $\Dmax'$ & \eqref{eq:printcurve} \\
Neutral axis warm/cool & $\delta_{RG}, \delta_{BG}$ & \eqref{eq:neutralaxis} \\
Silver retention & $\varrho$ & \eqref{eq:silverretention} \\
Grain amount & $g$ & \eqref{eq:grainapply} \\
Grain size & $r_g$ & \eqref{eq:grainkernel} \\
Halation threshold & $\tau$ & \eqref{eq:halsource} \\
Halation radius & $\ell$ & \eqref{eq:halpsf} \\
Halation intensity & $\alpha_h$ & \eqref{eq:haladd} \\
Halation color bias & $\boldsymbol{\beta}$ & \eqref{eq:haladd} \\
Diffusion & $w_d, \sigma_d$ & \eqref{eq:diffusion} \\
Vignetting & $v$ & \eqref{eq:vignette} \\
Chromatic aberration & $\varkappa$ & \eqref{eq:ca} \\
Barrel / pincushion & $k_1, k_2$ & \eqref{eq:distortion} \\
\bottomrule
\end{tabularx}
\end{center}
\end{table}
